% 01_intro.tex
\section{Introduction}
\label{sec:intro}

Quantum signal processing (QSP) applies a polynomial to a scalar
signal by interleaving a fixed single-qubit rotation, whose angle
encodes the signal, with single-qubit phase rotations whose angles
are free parameters~\cite{LowChuang2017,GSLW2019}.  The signal
$\tilde x$ lives on the interval $[-1,1]$, and after $d$ rotations
one matrix element of the circuit is a polynomial $P(\tilde x)$ of
degree at most $d$ with $|P|\le 1$.  Conversely, a completeness
theorem states exactly which polynomials, characterized by a
degree, a parity, and a normalization condition, are reachable by
some choice of phases.  In the
quantum singular value transformation the signal is an eigenvalue
or singular value of a block-encoded matrix, and the same
construction applies a polynomial to the whole
matrix~\cite{GSLW2019,MartynRossiTanChuang2021}.  The interval
$[-1,1]$ is therefore built into every QSP-based algorithm.  A
target defined on the whole real line, such as a Lorentzian or a
Gaussian amplitude profile, must first be truncated to a finite
window and rescaled into the interval, and the cost then grows
with the window size~\cite{McArdleGilyenBerta2025}.

It is natural to ask whether a change of variable removes this
restriction: whether a target that is costly as a polynomial in
$\tilde x$ becomes cheap when written in a coordinate that covers
the whole line.  Three coordinates for the same QSP circuit are
already in use.  The original formulation of Low and
Chuang~\cite{LowChuang2017} works with the rotation angle $\theta$
and produces trigonometric polynomials.  The formulation of
Gily\'en, Su, Low, and Wiebe~\cite{GSLW2019} works with
$\tilde x=\cos\theta$ and produces ordinary polynomials.
Generalized QSP~\cite{MotlaghWiebe2023} works with
$e^{i\theta}$ on the unit circle, allows arbitrary $\SU(2)$
rotations in place of the phase rotations, and produces Laurent
polynomials; this form was identified with the $\SU(2)$ nonlinear
Fourier transform by Alexis, Mnatsakanyan, and
Thiele~\cite{AlexisMnatsakanyanThiele2024,AlexisBeckerEtAl2026} and
by Laneve~\cite{Laneve2025b}, and its phases are computed by the
layer-stripping algorithm of Ni, Sarkar, Ying, and
Lin~\cite{NiSarkarYingLin2025}.  None of these coordinates covers
the real line.  The coordinate that does is $x=\cot\theta$, which
maps the circle onto $\R\cup\{\infty\}$ and in which the natural
basis is Boyd's family of rational Chebyshev
functions~\cite{Boyd1987b,Boyd1990}.  We call it the stereographic
coordinate, because $x$ is the stereographic projection of the
point $e^{2i\theta}$ of the unit circle onto the real line.
Figure~\ref{fig:charts} collects the four coordinates and the
changes of variable that relate them.

Here we answer the question in two parts.  First, we write QSP in
the stereographic coordinate and show that the resulting circuit
is the standard QSP circuit conjugated by one fixed gate.  Its
achievable polynomials, its phases, and its completeness theorem
therefore follow from the standard theory.  The circuit itself
remains bounded; what is unbounded is the ratio $P/Q$ of the two
amplitudes of its output qubit, which is a rational function of
$x$ on the whole real line and is formed after measurement at a
sampling cost that we quantify.  Second, we show that for
functions of a block-encoded matrix the change of variable gives
no advantage.  A polynomial of degree $d$ in $\tilde x$ is a rational
Chebyshev expansion of degree $d$ in $x$ with the \emph{same}
coefficients, so the achievable set and the approximation error at
every degree are identical in all four coordinates.

The freedom the change of variable does not touch is the choice of
sample points, and we show that state preparation can exploit it.
To load a function $f$ on $\R$ into the amplitudes of an $n$-qubit
register we sample it on the grid $y_j=L\cot\theta_j$, with
$\theta_j$ uniformly spaced, and weight each sample by the square
root of the local grid spacing.  The amplitudes are then uniform
samples of the function $G(\theta)=f(L\cot\theta)/\sin\theta$.
They are loaded either by a quantum Fourier transform of the
Fourier coefficients of $G$, or by a Fourier-type QSP circuit whose
signal is a diagonal phase operator costing $\mathcal{O}(n)$ gates
per call.  This operator is the Cayley transform of the position
operator of the grid, so the second circuit is QSP of an unbounded
operator.  In both circuits the cost is set by the smoothness of
$G$, which Boyd's theory of rational Chebyshev expansions
classifies.  The weight makes the map from $f$ to $G$ unitary, and
the grid sums that define norms and overlaps are exact, so the
state represents $f$ in the norm of $L^2(\R)$ with no
discretization error and no truncation radius.  The representation
is by interpolation in rational Chebyshev functions: the state
gives overlaps and a class of expectation values as exact
integrals over $\R$, but it is not a set of samples of $f$ on a
uniform grid.  For the Lorentzian the weighted function is
$\sin\theta$, and the state is exact at $\mathcal{O}(n)$ gates.
The success amplitude of the QSP circuit equals $M^{-1/2}$, where
$M$ is the constant of classical rejection sampling of $|f|^2$ from
a Cauchy proposal; it is of order one for a single peak with a
tail no heavier than $1/|y|$.  Constructions on a uniform grid
sample $f$ over a finite window and represent the same function up
to the discarded tail; Sec.~\ref{subsec:comparison} compares the
costs at equal $L^2$ error.

Section~\ref{sec:projective} reviews standard QSP, introduces the
stereographic coordinate, and derives its completeness theorem and
read-out.  Section~\ref{sec:limits} shows that a change of variable
cannot accelerate a block-encoded matrix function.
Section~\ref{sec:stateprep} develops the state-preparation
algorithm and compares it with existing methods.
Section~\ref{sec:discussion} summarizes the results and their
limitations.  The extension to complex signals and its relation to
the nonlinear Fourier transform are given in
Appendix~\ref{app:thm-gen}.

\begin{figure*}[t]
	\centering
	\begin{tikzpicture}[
		>=Stealth,
		every node/.style={font=\small},
		picnode/.style={
			rectangle, rounded corners=2pt, draw=black!60,
			thick, align=center, inner sep=5pt, minimum width=3.1cm
		},
		fourier/.style    ={picnode, fill=violet!10},
		chebyshev/.style  ={picnode, fill=blue!10},
		gqsp/.style       ={picnode, fill=orange!10},
		ratcheb/.style    ={picnode, fill=teal!10},
		edgelab/.style    ={font=\footnotesize, midway, sloped},
	]
	\node[fourier]   (F)  at (-4.0,  1.8) {%
		\textbf{Fourier} \\[1pt]
		$\theta \in (0,\pi)$ \\
		$\cos k\theta,\ \sin k\theta$};
	\node[chebyshev] (C)  at ( 4.0,  1.8) {%
		\textbf{Chebyshev} \\[1pt]
		$\tilde x \in (-1,1)$ \\
		$T_k(\tilde x),\ U_{k-1}(\tilde x)\sqrt{1-\tilde x^2}$};
	\node[gqsp]      (G)  at (-4.0, -1.8) {%
		\textbf{Unit circle} \\[1pt]
		$w \in \T$ \\
		$w^k+w^{-k},\ w^k-w^{-k}$};
	\node[ratcheb]   (R)  at ( 4.0, -1.8) {%
		\textbf{Stereographic} \\[1pt]
		$x \in \R$ \\
		$\TB_k(x),\ \SB_k(x)$};

	\draw[<->, thick] (F) -- (C)
	node[edgelab, above] {$\tilde x = \cos\theta$}
	node[edgelab, below] {$\theta = \arccos \tilde x$};

	\draw[<->, thick] (F) -- (G)
	node[edgelab, above] {$w = e^{2i\theta}$}
	node[edgelab, below] {$\theta = \tfrac12\arg w$};

	\draw[<->, thick] (C) -- (R)
	node[edgelab, above] {$x = \tilde x/\sqrt{1-\tilde x^2}$}
	node[edgelab, below] {$\tilde x = x/\sqrt{1+x^2}$};

	\draw[<->, thick] (G) -- (R)
	node[edgelab, above] {$x = i(w+1)/(w-1)$}
	node[edgelab, below] {$w = (x+i)/(x-i)$};

	\draw[<->, dashed, thick, gray] (F) -- (R)
	node[edgelab, above, gray] {$x = \cot\theta$};
	\end{tikzpicture}
	\caption{Four coordinates for the signal of one QSP circuit,
		with the changes of variable on the edges and the natural
		basis of each.  The rotation angle $\theta$ (Fourier) gives
		trigonometric polynomials; $\tilde x=\cos\theta$
		(Chebyshev) gives the ordinary polynomials of standard
		QSP~\cite{GSLW2019}; a point $w$ of the unit circle gives
		Laurent polynomials, as in generalized
		QSP~\cite{MotlaghWiebe2023}; and $x=\cot\theta$
		(stereographic) gives Boyd's rational Chebyshev functions
		$\TB_k(x)=\cos(k\,\mathrm{arccot}\,x)$ and
		$\SB_k(x)=\sin(k\,\mathrm{arccot}\,x)$, which are orthogonal
		on $\R$ with weight $1/(1+x^2)$~\cite{Boyd1990}.  Each
		basis function is the image of $\cos k\theta$ or
		$\sin k\theta$, so the same circuit is described in every
		box; only the stereographic coordinate covers the whole
		real line.  We use $w=e^{2i\theta}$ so that the bottom
		edge is the Cayley transform; generalized QSP uses
		$e^{i\theta}$, whose square is $w$.}
	\label{fig:charts}
\end{figure*}
